\chapter{Antecedentes}\label{ch:antecedentes}\label{ch:marco-teorico}

Este cap\'itulo re\'une el marco conceptual y el estado del arte que sustentan Prevex. Primero se presentan los antecedentes de dominio y de tecnolog\'{\i}a (marco te\'orico); luego se sistematizan soluciones comerciales y trabajos acad\'emicos recientes (estado del arte).

\section{Marco te\'orico}\label{sec:marco-teorico}

El marco te\'orico organiza, en primer lugar, el dominio de la prevenci\'on en aseguradoras de riesgos del trabajo y la caracterizaci\'on de los datos; a continuaci\'on, los modelos predictivos, la explicabilidad y la arquitectura de software.

\subsection{Gesti\'on de riesgos laborales y el paradigma de la Prevenci\'on 4.0}

Hist\'oricamente, la gesti\'on de la seguridad en el \'ambito laboral ha estado dominada por modelos de inspecci\'on y control reactivos. En estos esquemas, la actuaci\'on de las autoridades y de las aseguradoras de riesgos del trabajo se daba principalmente como respuesta a la ocurrencia de un siniestro o a la recepci\'on de una denuncia, con el fin de mitigar consecuencias legales y econ\'omicas. Sin embargo, los altos costos asociados a la siniestralidad y el impacto del ausentismo en la productividad empresarial han impulsado la necesidad de evolucionar hacia modelos proactivos y planificados.

En este contexto, el Estado nacional ha comenzado a promover una modernizaci\'on del sector. Este cambio de paradigma fue formalizado mediante la Resoluci\'on 48/2025 de la Superintendencia de Riesgos del Trabajo (SRT), publicada en el \emph{Bolet\'{\i}n Oficial de la Rep\'ublica Argentina} el 4 de noviembre de 2025. Dicha normativa establece los lineamientos de la ``Prevenci\'on 4.0'', incentivando a las aseguradoras a integrar tecnolog\'{\i}as de la informaci\'on, an\'alisis de grandes vol\'umenes de datos e Inteligencia Artificial para anticipar eventos no deseados \parencite{srt_res48_2025}.

La aplicaci\'on de este marco normativo habilita el desarrollo de sistemas predictivos que optimicen la asignaci\'on de recursos limitados, en particular de inspectores y preventores, focalizando los esfuerzos de auditor\'{\i}a en aquellos establecimientos con mayor probabilidad de sufrir un accidente en el corto plazo, en l\'{\i}nea con las nuevas orientaciones de la SRT en materia de prevenci\'on tecnol\'ogica \parencite{srt_prevencion40,srt_res48_2025}.

En t\'erminos operativos, la adopci\'on de modelos de inteligencia artificial resulta relevante para la gesti\'on de programas regulatorios vigentes como el Programa de Empleadores con Siniestralidad Elevada (PESE). El PESE tiene por objetivo realizar el seguimiento y control de empleadores con altos niveles de siniestralidad, estableciendo acciones espec\'{\i}ficas de prevenci\'on orientadas a reducir la accidentalidad y mejorar las condiciones de salud y seguridad en el trabajo \parencite{srt_pese}.

Una vez incluido en el PESE, el empleador y su aseguradora deben cumplir obligaciones espec\'{\i}ficas durante un per\'{\i}odo de dos a\~nos, entre las que se destacan la confecci\'on del Informe General del Empleador (IGE), la suscripci\'on de un Plan de Reducci\'on de la Siniestralidad (PRS) con plan de visitas, la actualizaci\'on del Relevamiento General de Riesgos Laborales (RGRL) y el cumplimiento del plan de fiscalizaci\'on por parte de la aseguradora \parencite{srt_pese}. En este escenario, la implementaci\'on de un modelo predictivo permite identificar de manera temprana aquellos establecimientos con trayectorias de mayor riesgo. De este modo, la aseguradora puede intervenir preventivamente antes de que la empresa alcance los umbrales de inclusi\'on en el PESE, reduciendo costos de cumplimiento regulatorio y, fundamentalmente, evitando la materializaci\'on de da\~nos a la salud de los trabajadores.

\subsection{Origen y caracterizaci\'on de los datos}\label{sec:datos}

El \emph{dataset} tabular del modelo proviene de una empresa desarrolladora del sistema core utilizado por las aseguradoras de riesgos del trabajo, que concentra en sus bases de datos el registro operativo de siniestros, p\'olizas y establecimientos empleadores. La extracci\'on cubre historia de accidentes de trabajo y variables asociadas (2022--2025) y se anonimiz\'o antes de cualquier uso anal\'{\i}tico. Es informaci\'on primaria, generada en el marco de la relaci\'on entre aseguradora y empleador, sin fuentes p\'ublicas masivas ni datos obtenidos mediante \emph{web scraping}.

El panel operativo es del orden de 35.500 establecimientos. Cada observaci\'on se construye a nivel de establecimiento y per\'{\i}odo (mes), con variables de estructura laboral (dotaci\'on, antig\"uedad, rotaci\'on), actividad econ\'omica (Clasificador Industrial Internacional Uniforme, CIIU), trayectoria de siniestralidad (frecuencia, severidad, tipo de prestaci\'on), diagn\'osticos, edades de la plantilla y antecedentes de fiscalizaci\'on o inclusi\'on en programas como el PESE. La variable objetivo expresa si el establecimiento registra al menos un accidente en los tres meses siguientes al per\'{\i}odo de decisi\'on (horizonte trimestral).

Previo al modelado se aplic\'o depuraci\'on de valores nulos y duplicados, estandarizaci\'on de formatos, codificaci\'on de variables categ\'oricas y construcci\'on de variables de tendencia (por ejemplo, siniestros acumulados en ventanas m\'oviles). La partici\'on es temporal: los \'arboles se ajustan con 2022--2023; 2024 fija el umbral de decisi\'on, sin reentrenar; el desempe\~no que se reporta es el holdout 2025, que el modelo no vio.

\subsection{Modelos predictivos y machine learning para clasificaci\'on}\label{sec:modelos-ml}

Para lograr la anticipaci\'on de siniestros, se recurre al Aprendizaje Autom\'atico (\emph{Machine Learning}), una rama de la Inteligencia Artificial que permite a los sistemas aprender e identificar patrones complejos a partir de datos hist\'oricos sin ser expl\'{\i}citamente programados \parencite{bishop2006}.

En el contexto de este proyecto, el problema se modela como una tarea de clasificaci\'on binaria supervisada, donde el objetivo es predecir si un establecimiento laboral sufrir\'a o no un accidente en un per\'{\i}odo determinado. En problemas de riesgo laboral, este enfoque resulta especialmente \'util cuando existen m\'ultiples variables interrelacionadas, como n\'ominas, antecedentes de siniestros, rubro de actividad, caracter\'{\i}sticas de la planta laboral, diagn\'osticos, edades de trabajadores, entre otras, cuya combinaci\'on no puede capturarse de manera adecuada mediante reglas fijas o umbrales simples.

La literatura especializada en aprendizaje autom\'atico sobre datos tabulares ha mostrado de forma consistente que los modelos basados en \'arboles de decisi\'on ensamblados suelen ofrecer un desempe\~no competitivo y en muchos contextos superior frente a arquitecturas de aprendizaje profundo, especialmente cuando el volumen de datos es moderado, las variables son heterog\'eneas y la interpretabilidad del modelo resulta relevante \parencite{grinsztajn2022}. Dentro de esta familia de algoritmos, XGBoost (\emph{Extreme Gradient Boosting}) ocupa un lugar destacado por su eficiencia computacional y su capacidad de modelar relaciones no lineales entre variables. El m\'etodo combina secuencialmente m\'ultiples \'arboles de decisi\'on d\'ebiles, de modo que cada nuevo \'arbol corrige los errores residuales del conjunto anterior, optimizando una funci\'on de p\'erdida regularizada que contribuye a controlar el sobreajuste (\emph{overfitting}) \parencite{chen2016}.

En dominios como la siniestralidad laboral, adem\'as, suele presentarse un desbalance marcado entre clases: los accidentes constituyen eventos relativamente infrecuentes frente al conjunto de establecimientos sin siniestros en un per\'{\i}odo dado. En este escenario, la m\'etrica de exactitud (\emph{accuracy}) puede resultar enga\~nosa, ya que un modelo que prediga sistem\'aticamente la clase mayoritaria obtendr\'{\i}a un porcentaje elevado de aciertos sin detectar los casos de riesgo relevantes \parencite{he2009}. Por ello, en problemas de detecci\'on de eventos cr\'{\i}ticos resulta m\'as adecuado priorizar m\'etricas orientadas a la clase positiva, como la sensibilidad (\emph{recall}), que mide la proporci\'on de casos de accidente correctamente identificados. Ello permite minimizar los falsos negativos, es decir, establecimientos con riesgo real que el sistema clasificar\'{\i}a err\'oneamente como seguros, un error particularmente grave en un contexto preventivo.

En la evaluaci\'on del prototipo tambi\'en se reportan m\'etricas adicionales. El \emph{PR-AUC} (\'area bajo la curva precisi\'on-recall) resume el equilibrio entre precisi\'on y recall a lo largo de umbrales, y es m\'as informativo que el ROC-AUC cuando la clase positiva es rara. El \emph{Brier score} mide el error cuadr\'atico medio entre la probabilidad predicha y el resultado binario observado (valores m\'as bajos indican mejor calibraci\'on). El umbral F2 es el punto de corte de decisi\'on que prioriza el recall respecto de la precisi\'on (el recall pesa el doble), alineado a un contexto preventivo donde omitir un establecimiento de riesgo es m\'as costoso que una visita extra.

En s\'{\i}ntesis, la combinaci\'on de clasificaci\'on supervisada, modelos de ensamble para datos tabulares y m\'etricas sensibles al desbalance de clases constituye la base conceptual que justifica el enfoque predictivo adoptado en este trabajo.

\subsection{Inteligencia artificial explicable (XAI) y valores SHAP}

Uno de los principales obst\'aculos para la adopci\'on de modelos predictivos de alta complejidad como los algoritmos basados en \'arboles en entornos empresariales y regulatorios es su car\'acter de ``caja negra''. Si bien estos modelos pueden alcanzar un elevado desempe\~no predictivo, no ofrecen de manera intr\'{\i}nseca una explicaci\'on comprensible de los factores que conducen a una determinada clasificaci\'on. En el \'ambito de la prevenci\'on de riesgos laborales, esta limitaci\'on resulta especialmente relevante: un inspector o preventor no puede fundar de manera operativa una auditor\'{\i}a o una intervenci\'on preventiva apoy\'andose \'unicamente en un puntaje o probabilidad de riesgo desprovisto de interpretaci\'on.

Frente a esta problem\'atica surge el campo de la Inteligencia Artificial Explicable (\emph{Explainable Artificial Intelligence}, XAI), orientado a desarrollar m\'etodos que hagan comprensibles, transparentes y auditables los resultados producidos por sistemas de aprendizaje autom\'atico \parencite{arrieta2020}. En t\'erminos conceptuales, la XAI no busca reemplazar el juicio experto humano, sino complementarlo: traduce salidas algor\'{\i}tmicas en informaci\'on interpretable que permita a los responsables de la toma de decisiones evaluar si una predicci\'on es razonable, identificar las variables m\'as influyentes y comunicar el fundamento de una acci\'on preventiva.

Para hacer comprensibles las predicciones del modelo XGBoost, se utilizan los valores SHAP, una t\'ecnica que indica cu\'anto influye cada variable en el resultado de un caso concreto (por ejemplo, la cantidad de empleados, el historial de siniestros o el rubro de actividad) \parencite{lundberg2017}. De este modo, en lugar de mostrar solo que un establecimiento tiene alto riesgo, el sistema puede explicar qu\'e factores empujaron esa clasificaci\'on, como un aumento reciente de siniestros o un tama\~no de plantilla elevado. Eso permite que inspectores y preventores de la aseguradora entiendan el motivo de la alerta, prioricen acciones correctivas y justifiquen una auditor\'{\i}a con argumentos concretos, no solo con un n\'umero.

As\'{\i}, la incorporaci\'on de XAI mediante valores SHAP no solo responde a una necesidad t\'ecnica de interpretabilidad, sino que constituye un requisito de usabilidad, trazabilidad y confianza en soluciones predictivas destinadas a entornos regulados, donde las decisiones deben poder explicarse ante equipos de prevenci\'on, auditor\'{\i}a y cumplimiento normativo \parencite{arrieta2020,lundberg2017}.

\subsection{Arquitectura cloud computing y patrones de dise\~no de software}

Un modelo predictivo solo aporta valor si puede ser usado en la operaci\'on diaria de la aseguradora. Sin embargo, por restricciones de seguridad y por la complejidad de los sistemas heredados (\emph{legacy}), como los ERP de las aseguradoras, no siempre resulta viable integrar el algoritmo directamente dentro de esas plataformas. Por ello, resulta adecuado desplegar la soluci\'on como un sistema independiente en la nube, accesible mediante una aplicaci\'on web que consuma los resultados del modelo sin modificar la infraestructura existente.

La computaci\'on en la nube consiste en acceder a recursos inform\'aticos (servidores, almacenamiento y redes) a trav\'es de Internet de forma flexible y bajo demanda, sin que la organizaci\'on deba administrar toda la infraestructura f\'{\i}sica \parencite{mell2011}. La nube p\'ublica permite alojar el sistema predictivo en servidores remotos; en este proyecto, mediante instancias de c\'omputo el\'astico Amazon EC2, lo que facilita el despliegue, el mantenimiento y la escalabilidad seg\'un la cantidad de usuarios y consultas.

En cuanto al dise\~no del software, se adopta una arquitectura cliente-servidor en lugar de una aplicaci\'on monol\'{\i}tica. Esto significa separar el sistema en capas con responsabilidades distintas que se comunican mediante API REST, un est\'andar que permite intercambiar datos entre el \emph{frontend} y el \emph{backend} a trav\'es de solicitudes HTTP \parencite{fielding2000}. La capa de \emph{backend}, desarrollada en .NET con C\#, se encarga de la l\'ogica de negocio, la persistencia en base de datos relacional (SQL Server) y la ejecuci\'on o exposici\'on de las predicciones del modelo XGBoost. La capa de \emph{frontend}, construida con React, ofrece la interfaz gr\'afica mediante la cual preventores e inspectores consultan los niveles de riesgo y las explicaciones asociadas.

Esta separaci\'on facilita que cada componente evolucione de forma independiente, simplifica el mantenimiento y permite entregar una herramienta segura e interactiva, que opera sobre datos previamente anonimizados y no requiere acceso directo a los sistemas internos de la aseguradora.

\section{Estado del arte}\label{ch:estado-del-arte}\label{sec:estado-del-arte}

Esta secci\'on sistematiza investigaciones, desarrollos y prototipos recientes vinculados a la predicci\'on de accidentabilidad laboral, con el fin de identificar limitaciones y fundamentar la brecha que Prevex busca resolver.

\subsection{Soluciones comerciales y de mercado vigentes}

En el \'ambito comercial y corporativo de las aseguradoras de riesgos del trabajo en la regi\'on, la adopci\'on de sistemas predictivos avanzados a\'un se encuentra en una fase incipiente. La mayor\'{\i}a de las soluciones de software denominadas \emph{Core de Seguros} operan bajo motores de reglas est\'aticas. Estos sistemas clasifican el riesgo de las empresas a partir de cruces de bases de datos simples: por ejemplo, si la empresa pertenece a un rubro de alto riesgo seg\'un el clasificador CIIU, es parte del programa PESE o si ya super\'o un l\'{\i}mite de siniestros denunciados.

Si bien estas herramientas automatizan el control regulatorio, carecen de capacidad predictiva real: no aprenden de la trayectoria de los datos ni pueden identificar combinaciones no lineales de factores que derivan en un accidente. Las plataformas comerciales gen\'ericas de Business Intelligence y an\'alisis predictivo (por ejemplo DataRobot o SAP Analytics Cloud) incorporan m\'odulos de Machine Learning, pero suelen requerir perfiles de Ciencia de Datos para operarlas. Esto las aleja del usuario final objetivo: el preventor o inspector de campo, quien necesita una herramienta interpretable y accionable de forma inmediata.

La Tabla~\ref{tab:benchmark-comercial} resume esa comparaci\'on de forma sistem\'atica.

\begin{table}[htbp]
	\centering
	\caption{Comparaci\'on de enfoques comerciales frente a Prevex. \fuentepropia}
	\label{tab:benchmark-comercial}
	\begin{tabular}{@{}p{0.22\textwidth}p{0.24\textwidth}p{0.24\textwidth}p{0.24\textwidth}@{}}
		\hline
		\textbf{Atributo} & \textbf{Core (reglas)} & \textbf{Plataformas ML gen\'ericas} & \textbf{Prevex} \\
		\hline
		Predicci\'on & Umbrales y reglas fijas & Modelos entrenables & Ranking trimestral XGBoost \\
		Explicabilidad & Baja (reglas) & Variable / t\'ecnica & SHAP en lenguaje operativo \\
		Usuario objetivo & Operaciones / IT & Cient\'{\i}fico de datos & Preventor / coordinador \\
		\hline
	\end{tabular}
\end{table}

Este diagn\'ostico es coherente con lo se\~nalado en la literatura regional \parencite{zamora2024}: incluso en el \'ambito acad\'emico el foco sigue puesto en la detecci\'on matem\'atica de patrones, sin traducirlos en interfaces operativas para los equipos de prevenci\'on. En consecuencia, existe una brecha entre las capacidades tecnol\'ogicas disponibles y las herramientas efectivamente utilizadas por los profesionales de seguridad y salud en el trabajo.

\subsection{Antecedentes acad\'emicos recientes}

Para relevar el estado de la investigaci\'on cient\'{\i}fica, se analizaron trabajos de los \'ultimos cinco a\~nos en repositorios acad\'emicos de alto nivel (arXiv, IEEE Xplore, ScienceDirect, PubMed/PMC) y publicaciones regionales. La tendencia global demuestra un inter\'es creciente en reemplazar los m\'etodos estad\'{\i}sticos tradicionales por modelos de Machine Learning para la salud ocupacional.

\subsubsection{Aplicaci\'on de modelos predictivos y limitaciones sectoriales}

Diversos autores han abordado la predicci\'on de accidentes utilizando algoritmos de ensamble y clasificaci\'on tradicional. Trabajos como el de \textcite{subauste2023} mediante Redes Bayesianas, o \textcite{mosquera2021} en la industria papelera, demuestran la viabilidad t\'ecnica de clasificar riesgos. Sin embargo, estos estudios frecuentemente reportan como limitaci\'on la dependencia a un \'unico sector econ\'omico o la dificultad de los modelos simples para generalizar sin caer en sobreajuste (\emph{overfitting}). \textcite{sarkar2022} representa un avance al aplicar Random Forest sobre 66.405 registros de OSHA abarcando 20 categor\'{\i}as industriales en los Estados Unidos, logrando \emph{accuracy} de 0,895 para predicci\'on de hospitalizaci\'on. No obstante, su modelo clasifica la gravedad de lesiones pasadas, sin predecir eventos futuros a nivel de establecimiento empleador.

\subsubsection{Deep learning versus datos tabulares}

Enfoques de vanguardia, como el propuesto por \textcite{yapi2025}, han intentado utilizar Redes Neuronales Recurrentes para predecir accidentes bas\'andose en series temporales de inspecciones previas. Aunque los resultados son prometedores y validan la importancia de las auditor\'{\i}as como variable predictiva, los autores coinciden en que el alt\'{\i}simo requerimiento computacional y la naturaleza de ``caja negra'' de estas redes dificultan su adopci\'on \'agil por parte de entidades aseguradoras. Este hallazgo es consistente con la evidencia m\'as amplia de la literatura: \textcite{grinsztajn2022} demostraron que los modelos basados en \'arboles de decisi\'on ensamblados superan sistem\'aticamente a las arquitecturas de aprendizaje profundo en datos tabulares heterog\'eneos, que son precisamente el tipo de datos que manejan las aseguradoras. Asimismo, \textcite{zermane2022} corroboraron que Random Forest supera a redes neuronales multicapa en conjuntos de datos de seguridad ocupacional, con \emph{accuracy} entre el 96\,\% y el 100\,\% en la predicci\'on de ca\'{\i}das fatales.

\subsubsection{El desaf\'{\i}o de la explicabilidad}

La literatura regional reciente subraya que el foco actual sigue puesto en la detecci\'on matem\'atica de patrones y tendencias \parencite{zamora2024}. Esto evidencia que la barrera principal para la implementaci\'on real no es la precisi\'on estad\'{\i}stica, sino la falta de interpretabilidad. Los modelos fallan en proveer a los equipos de auditor\'{\i}a las razones espec\'{\i}ficas para tomar acci\'on preventiva en campo.

Esta limitaci\'on no es exclusiva de los trabajos regionales. A nivel global, la comunidad cient\'{\i}fica ha reconocido que el desaf\'{\i}o de la adopci\'on de modelos predictivos en entornos operativos y regulados no reside en alcanzar mayor \emph{accuracy}, sino en traducir las salidas del modelo en decisiones comprensibles para el usuario no t\'ecnico \parencite{arrieta2020}. En el dominio de la seguridad laboral, esto es particularmente cr\'{\i}tico: un preventor o inspector necesita entender por qu\'e un establecimiento fue marcado como de alto riesgo, no solo que lo fue. Sin esa explicaci\'on, la predicci\'on no puede fundamentar una auditor\'{\i}a, no puede comunicarse al empleador y no puede sostenerse ante una instancia regulatoria como la SRT.

\subsection{Matriz de antecedentes}

A continuaci\'on se sistematizan siete enfoques representativos recientes, tanto a nivel internacional como regional, evaluando sus metodolog\'{\i}as, resultados y limitaciones con el fin de establecer la brecha de investigaci\'on que aborda el presente proyecto. La Tabla~\ref{tab:matriz-antecedentes} resume esa comparaci\'on.

\begin{landscape}
	\footnotesize
	\setlength{\tabcolsep}{3pt}
	\renewcommand{\arraystretch}{1.15}
	\begin{longtable}{@{}p{1.6cm}p{2.4cm}p{2.2cm}p{2.4cm}p{2cm}p{2.2cm}p{2.4cm}p{2.6cm}@{}}
		\caption{Matriz de antecedentes}\label{tab:matriz-antecedentes} \\
		\hline
		\textbf{Autor y A\~no} & \textbf{T\'{\i}tulo} & \textbf{Objetivo del Trabajo} & \textbf{Enfoque / Metodolog\'{\i}a} & \textbf{Tecnolog\'{\i}as Utilizadas} & \textbf{Resultados Principales} & \textbf{Limitaciones Reconocidas} & \textbf{Aporte para el Proyecto Propio} \\
		\hline
		\endfirsthead
		\caption[]{Matriz de antecedentes (continuaci\'on)} \\
		\hline
		\textbf{Autor y A\~no} & \textbf{T\'{\i}tulo} & \textbf{Objetivo del Trabajo} & \textbf{Enfoque / Metodolog\'{\i}a} & \textbf{Tecnolog\'{\i}as Utilizadas} & \textbf{Resultados Principales} & \textbf{Limitaciones Reconocidas} & \textbf{Aporte para el Proyecto Propio} \\
		\hline
		\endhead
		\hline
		\endfoot
		Yapi, A. (2025) arXiv &
		A new machine learning framework for occupational accidents forecasting with safety inspections integration &
		Proponer un marco predictivo de accidentes laborales integrando como variable la frecuencia y resultados de las inspecciones de seguridad. &
		Desarrollo de un modelo predictivo con datos hist\'oricos de inspecciones y siniestros (Francia). Validaci\'on con m\'etricas de clasificaci\'on binaria. &
		Redes Neuronales Recurrentes (LSTM), modelos basados en \'arboles (Random Forest, XGBoost), Python. &
		Los modelos de \'arbol superan a LSTM en viabilidad operativa. Incluir inspecciones previas mejora la predicci\'on hasta un 18\,\% de ganancia en \emph{recall}. &
		Las redes profundas exigen gran capacidad de c\'omputo, operan como ``caja negra'' y fueron validadas bajo marcos regulatorios extranjeros (Francia). &
		Valida la inclusi\'on de auditor\'{\i}as como variable predictiva. Refuerza la brecha: necesidad de aplicar SHAP y adaptar el modelo al marco normativo argentino (SRT/PESE). \\
		\hline
		Zamora-Cevallos, E. (2024) Revista Metropolitana de Ciencias Aplicadas &
		Modelo predictivo para determinar patrones y tendencias que influyen en los accidentes de trabajo &
		Identificar mediante datos hist\'oricos los patrones y tendencias principales que derivan en siniestralidad laboral en Ecuador. &
		An\'alisis de datos hist\'oricos de accidentabilidad. Enfoque cuantitativo con modelos de Machine Learning y regresi\'on lineal. &
		Modelos de Machine Learning (clasificadores), Regresi\'on Lineal. &
		Identificaci\'on de patrones estad\'{\i}sticos en la siniestralidad. Detecci\'on de variables de mayor incidencia. &
		El enfoque recae netamente en las previsiones matem\'aticas, careciendo de interfaces o herramientas para explicar las decisiones al personal de campo. &
		Evidencia que el sector se centra en detectar patrones matem\'aticos. Subraya el diferencial del presente proyecto: la incorporaci\'on de XAI (valores SHAP) para traducir esos patrones a los prevencionistas. \\
		\hline
		Subauste Oliden, D. A. (2023) UPC (Tesis de grado) &
		Sistema de predicci\'on de riesgos de accidentes laborales en el sector manufactura utilizando t\'ecnicas de redes bayesianas y \'arbol de decisi\'on &
		Desarrollar un sistema de predicci\'on de riesgos laborales enfocado exclusivamente en el sector manufacturero del Per\'u. &
		Dise\~no y entrenamiento de modelos de clasificaci\'on sobre datos sectoriales de manufactura. Evaluaci\'on con m\'etricas est\'andar. &
		Redes Bayesianas, \'Arboles de Decisi\'on simples. &
		Viabilidad t\'ecnica de clasificar riesgos en el sector manufactura. Precisi\'on limitada por el uso de algoritmos simples. &
		La precisi\'on predictiva se acota al utilizar algoritmos de clasificaci\'on tradicionales, que presentan dificultades con el sobreajuste (\emph{overfitting}) en datos heterog\'eneos. &
		Justifica la elecci\'on de XGBoost como algoritmo de ensamble superador que maneja el sobreajuste eficientemente en datos corporativos heterog\'eneos. \\
		\hline
		Sarkar, S.; Maiti, J. (2022) IJERPH (MDPI) &
		Occupational Injury Risk Mitigation: Machine Learning Approach and Feature Optimization for Smart Workplace Surveillance &
		Proponer un modelo predictivo de gravedad de lesiones laborales multisector con optimizaci\'on de variables, utilizando 66\,405 registros de OSHA (EE.UU.). &
		An\'alisis multisector con datos de OSHA (20 categor\'{\i}as NAICS). Comparaci\'on de cinco algoritmos. Optimizaci\'on de hiperpar\'ametros. &
		Support Vector Machine, K-Nearest Neighbors, Na\"{\i}ve Bayes, \'Arbol de Decisi\'on, Random Forest. Dataset OSHA (EE.UU.). &
		Random Forest super\'o a los dem\'as modelos (\emph{accuracy} 0,895 hospitalizaci\'on; 0,954 amputaci\'on). Las 3 variables m\'as influyentes: naturaleza de la lesi\'on, tipo de evento, parte del cuerpo afectada. &
		El modelo se entrena con datos estructurados de reportes post-accidente, no predice eventos futuros a nivel de establecimiento. Aplica la legislaci\'on norteamericana (OSHA). &
		Primer antecedente directo de aplicaci\'on multisector con Random Forest. Confirma la superioridad de los m\'etodos de ensamble. Diferencia: el presente proyecto predice accidentabilidad futura a nivel de empleador, no clasifica lesiones pasadas. \\
		\hline
		Mosquera Navarro, R. (2021) Informaci\'on Tecnol\'ogica &
		Predicci\'on de la accidentalidad laboral en la industria de pulpa y papel usando algoritmos de clasificaci\'on &
		Predecir accidentes acotando el modelo a las variables espec\'{\i}ficas de la industria de pulpa y papel en Colombia. &
		An\'alisis de datos hist\'oricos de una sola planta industrial. Entrenamiento y evaluaci\'on de clasificadores tradicionales. &
		Algoritmos de clasificaci\'on tradicionales y anal\'{\i}tica de datos. &
		Viabilidad de aplicar ML a la predicci\'on de accidentes en el sector papelero. Alta dependencia del entorno de una \'unica planta. &
		Alta dependencia de las variables de entorno de una sola planta, lo que reduce la escalabilidad del sistema hacia otros rubros. &
		Demuestra la viabilidad por sector. Resalta que el presente sistema aborda la brecha multi-rubro (cartera CIIU completa de una ART). \\
		\hline
		Su\'arez S\'anchez, F. A. (2021) Universidad de Medell\'{\i}n &
		Modelo de simulaci\'on y predicci\'on de riesgos laborales en la construcci\'on de edificios: caso Colombia &
		Predecir y simular riesgos laborales en obras de construcci\'on en Colombia mediante simulaci\'on estad\'{\i}stica. &
		Simulaci\'on de Montecarlo, tablas de contingencia (Chi-cuadrado), encuestas de campo. &
		M\'etodo de Montecarlo, software estad\'{\i}stico SPSS. &
		Identificaci\'on de factores de riesgo predominantes en el sector construcci\'on colombiano. Modelo de simulaci\'on funcional para ese contexto. &
		El modelo est\'a fuertemente anclado a un \'unico sector y a las variables demogr\'aficas y normativas de Colombia. No aplica ML. &
		Confirma la utilidad de la predicci\'on en la regi\'on. Diferencia metodol\'ogica clara: evoluci\'on de la simulaci\'on estad\'{\i}stica sectorial hacia un modelo de ML (XGBoost) escalable y multisector. \\
		\hline
		Zermane, A. et al. (2022) Safety Science (Elsevier) &
		Predicting fatal fall from heights accidents using random forest classification machine learning model &
		Predecir ca\'{\i}das fatales desde alturas en la industria de la construcci\'on en Malasia utilizando un modelo de Random Forest. &
		An\'alisis de registros del Departamento de Salud y Seguridad Ocupacional de Malasia. Comparaci\'on de algoritmos de ensamble vs.\ \'arboles simples. &
		Random Forest, MLP (redes neuronales), \'Arboles de Decisi\'on. Dataset nacional de Malasia. &
		Random Forest super\'o a MLP y \'Arboles de Decisi\'on, logrando \emph{accuracy} entre 96\,\% y 100\,\%. Los factores m\'as influyentes fueron edad, posici\'on del trabajo y experiencia. &
		Modelo exclusivo del sector construcci\'on y de la legislaci\'on malaya. Sin componente de explicabilidad (XAI). Resultados no generalizables a otros rubros. &
		Refuerza la superioridad de Random Forest y XGBoost sobre redes neuronales para datos tabulares sectoriales. Confirma la necesidad de adaptar el modelo a la normativa local (en este caso, SRT Argentina). \\
		\hline
	\end{longtable}
\end{landscape}


\subsection{Conclusiones del estado del arte}

La revisi\'on de la literatura y del mercado actual permite extraer cinco conclusiones fundamentales que justifican el alcance y el valor diferencial de este Proyecto Final de Ingenier\'{\i}a:

\begin{enumerate}
	\item \textbf{Validaci\'on del enfoque t\'ecnico:}
	Existe evidencia acad\'emica convergente de que las arquitecturas de ensamble son m\'as estables que los \'arboles de decisi\'on simples \parencite{subauste2023}, m\'as viables operativamente que las redes neuronales profundas \parencite{yapi2025} y que Random Forest y XGBoost superan a otras familias de algoritmos incluso en conjuntos de datos multisectoriales de gran volumen \parencite{sarkar2022,zermane2022}. Esto ratifica la elecci\'on de XGBoost como el motor predictivo \'optimo para datos tabulares corporativos heterog\'eneos.

	\item \textbf{Brecha de explicabilidad (XAI):}
	Las soluciones actuales, tanto comerciales como acad\'emicas, maximizan la precisi\'on estad\'{\i}stica pero descuidan la experiencia del usuario y la trazabilidad de las decisiones \parencite{zamora2024}. La incorporaci\'on de valores SHAP en el presente proyecto resulta indispensable para romper la caja negra y dotar a los inspectores de fundamentos auditables, en l\'{\i}nea con las demandas del campo de la Inteligencia Artificial Explicable \parencite{arrieta2020}.

	\item \textbf{Brecha regulatoria y de escala:}
	Gran parte de los modelos vigentes est\'an limitados a nichos industriales espec\'{\i}ficos \parencite{mosquera2021,zermane2022} o responden a legislaciones extranjeras \parencite{yapi2025,sarkar2022}. Este proyecto aporta un enfoque escalable a m\'ultiples rubros (CIIU) y parametrizado exclusivamente para la normativa argentina, incluyendo la Resoluci\'on SRT 48/2025 y el programa PESE.

	\item \textbf{Brecha de arquitectura e implementaci\'on:}
	La mayor\'{\i}a de las investigaciones culminan en documentos experimentales aislados o estudios te\'oricos sin implementaci\'on productiva. El presente proyecto llena esta brecha encapsulando el modelo anal\'{\i}tico dentro de una arquitectura Cliente-Servidor robusta y desacoplada (utilizando .NET para el \emph{backend} y React para el \emph{frontend}), lista para integrarse \'agilmente en el ecosistema real de una aseguradora.

	\item \textbf{Tendencia emergente de predicci\'on a nivel de empleador:}
	Ninguno de los trabajos relevados aborda la predicci\'on prospectiva de accidentabilidad a nivel de establecimiento empleador en el contexto de una aseguradora de riesgos del trabajo. Los proyectos existentes se orientan a clasificar la severidad de accidentes ya ocurridos, no a anticipar qu\'e empleadores de una cartera aseguradora presentar\'an siniestros en el pr\'oximo per\'{\i}odo. Esta es la contribuci\'on central y diferencial del presente proyecto.
\end{enumerate}
